% ECONOMICS_PROBLEM: {"id": "D83-640", "title": "AI-generated marketplace summaries", "jel_code": "D83", "source_id": "OP-0403"}
% FIRST_POSTED: 2026-10-09
% ATTEMPT_STATUS: unresolved
% ENTRY_KIND: research_attempt
\documentclass[11pt]{article}
\usepackage[T1]{fontenc}
\usepackage{amsmath,amssymb,amsthm}
\usepackage[margin=1in]{geometry}
\usepackage{hyperref}
\newtheorem{proposition}{Proposition}
\newtheorem{lemma}{Lemma}
\newcommand{\E}{\mathbb{E}}
\newcommand{\Pp}{\mathbb{P}}
\newcommand{\Var}{\operatorname{Var}}
\newcommand{\Cov}{\operatorname{Cov}}
\title{AI-generated marketplace summaries\\\large Information value and interference in marketplace summary experiments}
\author{GPT-6 (Codex)}
\date{October 9, 2026}
\begin{document}
\maketitle

\section{Policy, cohort, and welfare target}
The target is the six-month consumer-welfare effect of a specified
AI-summary policy, including users who do not buy. Fix eligibility before
the summary is displayed, and retain everyone in that cohort. Freeze the
model version, prompt, review cutoff, and refresh rule. For user $i$, let
$Y_i$ be independently assessed money-valued satisfaction less price and
search cost, including returns and disappointing purchases. The no-purchase
alternative has zero gross consumption value and price but still incurs
any search cost. This attempt establishes a rational-information benchmark,
an explicit failure example, and an estimator under an interference model.
It does not determine the actual six-month effect or seller response.

\section{A benchmark in which optional summaries cannot reduce welfare}
Consider a fixed review corpus, fixed offers, a finite state space $\Omega$,
a finite action set including no purchase, and utility $u(a,\omega)$.
Let $R$ denote the complete review information and let a summary $S$ be
produced by a stochastic kernel depending only on $R$. With correct beliefs
and no processing costs, write
\[
 V(R)=\E\max_a\E[u(a,\omega)\mid R],\qquad
 V(S)=\E\max_a\E[u(a,\omega)\mid S].
\]
Since $\omega,R,S$ form a Markov chain, conditional averaging implies
\[
 V(R)\ge V(S).                                           \tag{1}
\]
To prove this, choose an $S$-optimal action for each signal. An observer of
$R$ can simulate $S$ from its known kernel and take that action. This
achieves $V(S)$, while choosing an optimal action given $R$ does at least
as well. Equivalently, apply convexity of the maximum to posterior expected
utilities. The summary need not add statistical information to add value
when acquiring or interpreting the underlying reviews is costly.

A stronger conclusion applies to the catalogue's optional-access design.
Let $\mathcal P_0$ be the feasible information-acquisition and action
policies under the old interface, with their full search costs. If every
old policy remains available at the same cost after summaries are added,
then $\mathcal P_0\subseteq\mathcal P_1$, so
\[
 \sup_{\pi\in\mathcal P_1}\E[u(\pi,\omega)-c(\pi)]
 \ge\sup_{\pi\in\mathcal P_0}\E[u(\pi,\omega)-c(\pi)].     \tag{2}
\]
No informational ordering between summaries is needed for (2). Its proof
is set inclusion, and its assumptions are demanding: correct decision
models, free ability to ignore summaries, unchanged old-interface costs,
and no price or seller-quality response. Keeping a link to all reviews
does not establish those behavioral and institutional conditions.

\section{A checked negative-welfare example}
Let product quality be high with probability $1/4$ and low otherwise.
Gross values are three and zero, and the price is one, giving purchase
surpluses two and minus one. A full review perfectly reveals quality and
costs $0.1$ to read. Under the control interface, reading and buying only
high-quality products gives expected welfare
\[
 (1/4)2-0.1=0.4,
\]
which exceeds both no search/no purchase, with payoff zero, and uninformed
purchase, with payoff $-0.25$. The optimal purchase rate is $1/4$.

Suppose a misleading summary causes consumers to believe both that the
high-quality probability is $0.8$ and that the full review contains no
information beyond the summary. Their perceived immediate purchase
surplus is $0.8(2)+0.2(-1)=1.4$. Under these stated beliefs, further
inspection leaves that perceived surplus unchanged and adds cost $0.1$,
so buying immediately is subjectively optimal. Actual mean welfare becomes
\[
 (1/4)2+(3/4)(-1)=-0.25,
\]
a decrease of $0.65$, while purchases rise from $1/4$ to one. Full reviews
remain available throughout. This example intentionally violates correct
beliefs in (2), through both quality beliefs and beliefs about incremental
information. It shows why greater conversion and nominal review access
are insufficient welfare evidence. It does not claim that all AI summaries
produce these errors.

The control cohort's nonbuyers have paid the review cost. Dropping them
would replace $0.4$ by a selected purchaser average and conceal the proper
policy comparison. Audits should therefore measure omitted adverse
information and consumer interpretation as well as literal falsehoods.
A technically true summary can change beliefs about what further search
would reveal.

\section{A design-based result allowing cross-cluster exposure}
Let clusters $j=1,\ldots,J$ be independently assigned $Z_j\sim
\operatorname{Bernoulli}(p)$, with $0<p<1$. For each baseline user fix a
nonempty set $N_i$ of potentially relevant clusters before assignment.
Assume the six-month potential outcome is $Y_i(z_{N_i})$: all interference
relevant to that user is confined to this set. Define
\[
 E_i^1=\{Z_j=1\ \forall j\in N_i\},\quad
 E_i^0=\{Z_j=0\ \forall j\in N_i\},
\]
with probabilities $\pi_i^1=p^{|N_i|}$ and
$\pi_i^0=(1-p)^{|N_i|}$. The finite-cohort contrast between full exposure
to summaries and no exposure is
\[
 \tau_N=\frac1n\sum_i[Y_i(\mathbf1_{N_i})-Y_i(\mathbf0_{N_i})].
\]
Then the Horvitz--Thompson estimator
\[
 \widehat\tau_N=\frac1n\sum_i\left\{
 \frac{\mathbf1(E_i^1)Y_i^{\rm obs}}{\pi_i^1}
 -\frac{\mathbf1(E_i^0)Y_i^{\rm obs}}{\pi_i^0}\right\}       \tag{3}
\]
is unbiased. On $E_i^a$, consistency and the exposure restriction make the
observed outcome equal the corresponding fixed potential outcome;
taking expectations cancels the exposure probability in each summand.
No independence between users is needed for this expectation calculation.

Exposure overlap can nevertheless make precision poor. Under the extra
bound $|Y_i(z)|\le B$, define
\[
 c_j=\frac{B}{n}\sum_{i:j\in N_i}
       \left(\frac1{\pi_i^1}+\frac1{\pi_i^0}\right).
\]
Changing one assignment $Z_j$ can change only summands for users with
$j\in N_i$. Each such summand, before division by $n$, takes one of the
three values $Y_i(\mathbf1)/\pi_i^1$, zero, or
$-Y_i(\mathbf0)/\pi_i^0$; their pairwise distances are bounded by the
quantity defining $c_j$. Bounded differences for independent cluster
assignments therefore gives
\[
 \Pp(|\widehat\tau_N-\tau_N|\ge t)
 \le2\exp\left(-\frac{2t^2}{\sum_jc_j^2}\right).           \tag{4}
\]
This is conservative but fully design-based under the stated assumptions.
For one cluster per user, weights are ordinary inverse assignment
probabilities. For users linked to many clusters, the exponential growth
of inverse exposure probabilities explains why more observations may
still leave little useful precision.

The set $N_i$ must include potential exposure routes, not merely products
a user actually viewed after treatment. Summary-induced search can change
those realized sets. Selecting $N_i$ after observing views would invalidate
the assignment-probability calculation and could condition on a treatment
mediator. Likewise, market-wide prices or sellers influenced by clusters
outside $N_i$ invalidate the assumed outcome mapping.

\section{What remains for the original problem}
A frozen-corpus trial estimates a direct interface effect only under its
fixed offers and review supply. A market-level six-month trial can include
seller adaptation, review manipulation, and changing availability, but its
assignment unit must encompass the relevant feedback. One cannot transfer
(3) to that ecosystem without expanding its exposure sets or defending a
new interference restriction. Broad exposure sets may make (4) vacuous;
that is an identification-and-design limitation, not a reason to silently
assume individual randomization.

Farronato's 2025 conference paper, Section 1.1, explicitly identifies
AI aggregation of textual reviews as a research direction and discusses
average/distributional effects and interference. That source motivates
this problem; it does not supply the welfare sign. Equations (1)--(4)
are restricted decision-theoretic and randomization calculations.
Actual monetary satisfaction, search costs, misleading omissions, group
heterogeneity, and seller adaptation remain to be measured under a
specified policy. The original empirical welfare question is unresolved.

\section{Completion Estimate}
40\%

This is a post hoc, heuristic estimate for the full catalogue target.
The attempt proves optional-information welfare benchmarks, constructs a misleading-summary welfare loss and gives unbiased full-exposure estimation with design-based interference bounds. Actual immediate and six-month group welfare, factual omissions and seller or review adaptation still require measured outcomes and exposure designs spanning the relevant marketplace feedback.

\begin{thebibliography}{9}
\bibitem{f} Chiara Farronato (2025),
\emph{Digital Platforms as Information Aggregators: Implications for their
Design and Regulation}, conference-paper version, Section 1.1.
\url{https://conference.nber.org/confer/2025/DTs25/farronato.pdf}.
\end{thebibliography}

\end{document}
